% ============================================================================
%  Chapter 2 -- First Contact: RISC-V Assembly on PRG32          (Lecture 02)
% ============================================================================
\chapter{First Contact: RISC-V Assembly on PRG32}
\label{ch:firstcontact}

\epigraphline{Talk is cheap. Show me the code.}
  {Linus Torvalds, Linux kernel mailing list, 2000}

\begin{chapgoals}
By the end of this chapter you will be able to
\begin{itemize}
\item build a cartridge from an assembly file and run it on \prg;
\item say what a register is, and name the ones used every day;
\item call a function of the framework, passing arguments in registers;
\item keep a value in memory from one frame to the next: load, modify, store;
\item read one instruction as the 32-bit number that it really is;
\item list the questions that this program raises, which the rest of the book answers.
\end{itemize}
\end{chapgoals}
% Generated by tools/textbook/export_listings.js from tools/decks/mastery.js and tools/concepts/registry.json: do not edit.
\begin{chapstatus}{Learning status of this chapter}
\masterychip{now}{UNDERSTAND NOW}\enspace registers, compute, move, jump; registers; functions, call and return; the init, update, draw contract; reading a line of assembly; arguments in a0 to a7; why ra is saved; address, value, load, store; the base instructions; pseudo-instructions; an instruction is a number; tracing by hand.\par\smallskip
\masterychip{use}{USE NOW}\enspace hexadecimal (Chapter~\ref{ch:information}); bit, byte, word (Chapter~\ref{ch:information}); the calling convention (Chapter~\ref{ch:processor}); bit masks (Chapter~\ref{ch:boole1}) and, used without being studied in this course, the assembler directives, the PRG32 tools, the PRG32 runtime calls, the PRG32-QT debugger, the complete program.\par\smallskip
\masterychip{later}{UNDERSTAND LATER}\enspace the stack (Chapter~\ref{ch:processor}); addressing modes (Chapter~\ref{ch:instructions}); instruction formats (Chapter~\ref{ch:instructions}); two's complement (Chapter~\ref{ch:integers}); fixed point (Chapter~\ref{ch:reals}); Boolean algebra (Chapter~\ref{ch:boole1}); the adder (Chapter~\ref{ch:boole2}); from volts to bits (Chapter~\ref{ch:electronics}); latches and flip-flops (Chapter~\ref{ch:sequential}); the processor cycle (Chapter~\ref{ch:instructions}); the Verm1 project (Chapter~\ref{ch:verm1p1}).
\end{chapstatus}


This chapter is the ``play'' turn of the spiral for the whole first half of
the book: every later chapter points back to something typed here. The reader
is \emph{not} expected to understand everything in it. This chapter is for
seeing; the following ones are for explaining. Four small programs are
written, each complete and each running; they are the verified steps of the
first assembly laboratory of the \prg textbook \citep{prg32textbook}.

% ----------------------------------------------------------------------------
\section{The machine}

\subsection{Binary and hexadecimal: how to read the numbers}

Before the first program, a word on how its numbers are written, because
assembly code is full of them. A computer stores everything as \term{bits},
digits that are 0 or~1; a number written with bits is in \term{binary}, base~2,
where each position is worth twice the one on its right: $1010_2 = 8 + 2 = 10$.
Long rows of bits are hard to read, so programmers group them four at a time
and write each group as one \term{hexadecimal} digit, base~16: 0 to 9, then A to
F for ten to fifteen. The prefix \texttt{0x} says that a number is hexadecimal,
so \hex{10} is sixteen and not ten.

\begin{center}\small
\begin{tabular}{rll}
\toprule
\textbf{Decimal} & \textbf{Binary} & \textbf{Hexadecimal}\\
\midrule
0 & \texttt{0000} & \hex{0}\\
1 & \texttt{0001} & \hex{1}\\
9 & \texttt{1001} & \hex{9}\\
10 & \texttt{1010} & \hex{A}\\
15 & \texttt{1111} & \hex{F}\\
16 & \texttt{1 0000} & \hex{10}\\
255 & \texttt{1111 1111} & \hex{FF}\\
\bottomrule
\end{tabular}
\end{center}

\noindent One hexadecimal digit is exactly four bits: \hex{07E0}, the green of
\prg, is \texttt{0000 0111 1110 0000}. Today it is enough to recognise these
notations; Chapter~\ref{ch:information} explains them, with conversions in
both directions.

\subsection{Registers, and three kinds of instruction}

A \riscv processor, seen by the programmer, is a very small thing.

\begin{description}[style=nextline,leftmargin=1.5em]
\item[Thirty-two registers.] Each \term{register} holds one 32-bit value. The
  processor can compute only on values that are in registers.
\item[Compute.] An arithmetic or logic operation between registers:
  \ins{add}, \ins{sub}, \ins{and}, \ins{or}, \ins{xor}, and their
  \term{immediate} forms, such as \ins{addi}, in which one operand is a small
  constant written in the instruction.
\item[Move.] A transfer between a register and memory: \ins{lw} loads a word,
  \ins{sw} stores a word. Nothing else touches memory.
\item[Jump.] A change in which instruction runs next: branches such as
  \ins{beqz}, calls and returns.
\end{description}

That is very nearly the whole machine. The richness of computing comes not
from a large instruction set but from combining a small one in long sequences.

\subsection{The registers actually used}

The hardware calls the registers \reg{x0} to \reg{x31} and, with one
exception, treats them alike. Programmers call them by other names, which
record the \emph{role} that each plays (Table~\ref{tab:fc:regs}). The roles
are fixed not by the hardware but by an agreement, the \term{calling
convention}, part of the application binary interface (ABI). Following it is
what allows a function written in assembly and the C code of the framework to
call each other.

\begin{table}[htb]
\centering\small
\caption{The registers used every day. The complete table is in
Appendix~\ref{app:riscv}.}\label{tab:fc:regs}
\begin{tabularx}{\textwidth}{@{}llLl@{}}
\toprule
\textbf{Register} & \textbf{Name} & \textbf{Role} & \textbf{Preserved across a call?}\\
\midrule
\reg{a0}--\reg{a7} & arguments & values passed into a call; \reg{a0} also carries the return value & no\\
\reg{t0}--\reg{t6} & temporaries & scratch values & no\\
\reg{s0}--\reg{s11} & saved & values that a call promises to preserve & yes\\
\reg{ra} & return address & where \ins{ret} jumps back to & no: save it yourself\\
\reg{sp} & stack pointer & top of the stack, kept 16-byte aligned & yes\\
\reg{zero} & zero & always reads as 0; writes are ignored & --\\
\bottomrule
\end{tabularx}
\end{table}

\subsection{Functions: a piece of program with a name}

A \term{function} is a sequence of instructions that does one job, such as
clearing the screen, and starts at a label, which gives it its name. To
\term{call} a function is to jump to its first instruction after noting where
to come back: the instruction \ins{call} does both. When its job is done the
function \term{returns}: \ins{ret} jumps back to the instruction that follows
the call. The \term{arguments} are the values a function needs for its job; in
\riscv they travel in the registers \reg{a0} to \reg{a7}. The code that calls
is the \term{caller}, the function called is the \term{callee}, and the place
to come back to is kept in the register \reg{ra}, the return address: this is
why step~2 saves \reg{ra} before its own calls and restores it before \ins{ret}.
\prg offers ready-made functions such as \code{prg32\_gfx\_clear}; you write
three of your own, which the runtime calls.

\subsection{The contract: three functions}

A \prg game is three functions that the programmer writes and the firmware
calls (Figure~\ref{fig:fc:contract}). \code{init} is called once, when the
cartridge starts, and sets up the state of the game. \code{update} is called
once per frame: it reads the input and changes the state. \code{draw} is
called once per frame: it turns the state into a picture. The firmware finds
the three functions by name: a prefix chosen by the programmer, followed by
\code{\_init}, \code{\_update} and \code{\_draw}. That is the whole contract
between a game and the machine.

\begin{figure}[htb]
\centering
\begin{tikzpicture}[node distance=10mm]
  \node[blkacc,text width=30mm,minimum height=14mm] (i) {\texttt{init}\\[1pt]\footnotesize once: set up the state};
  \node[blkamb,text width=30mm,minimum height=14mm,right=of i] (u) {\texttt{update}\\[1pt]\footnotesize every frame: input, state};
  \node[blkvio,text width=30mm,minimum height=14mm,right=of u] (d) {\texttt{draw}\\[1pt]\footnotesize every frame: state to picture};
  \draw[->,thick] (i) -- (u); \draw[->,thick] (u) -- (d);
  \draw[->,thick,rounded corners=3mm] (d.south) -- ++(0,-6mm) -| node[pos=0.25,below,lbl]{next frame} (u.south);
\end{tikzpicture}
\caption{The contract between a cartridge and the runtime.}
\label{fig:fc:contract}
\end{figure}

% ----------------------------------------------------------------------------
\section{Your first cartridge}

\subsection{How to read one line of assembly}

Before the first listing, a guide to reading it. A line of assembly has at most
four parts:

\begin{asm}
first_update:   addi t1, t1, 1      # one more tick
\end{asm}

\noindent A \term{label}, \code{first\_update:}, is a name for this place in
the program and ends with a colon. The \term{instruction}, \ins{addi}, is the
operation to perform, here ``add a constant''. The \term{operands},
\code{t1, t1, 1}, are what it works on: the destination first, then the
sources, so the line reads ``put \reg{t1} plus 1 into \reg{t1}'', the
assignment $t1 = t1 + 1$. A \term{comment} runs from \code{\#} to the end of the
line and is ignored by the assembler. A line that begins with a dot, such as
\code{.section .text}, is a \term{directive}: an order to the assembler, not to
the processor. Blank lines and indentation matter to the reader, not to the
assembler.

\subsection{Step 1: the smallest program that runs}

Listing~\ref{lst:fc:step1} is a complete cartridge. Create a folder, type the
program into a file called \fname{step1\_skeleton.S}, and resist the
temptation to copy it: the fingers learn too.

\asmfile[caption={Step 1: three functions that do nothing, correctly.},label=lst:fc:step1]{code/l02/c03.S}

Lines that begin with a dot are \term{directives}: orders to the assembler,
not instructions for the processor. \code{.option norelax} keeps the layout
of the cartridge predictable, and every \prg example starts with it.
\code{.section .text} says that what follows is executable code.
\code{.global} makes a \term{label} visible outside the file, so that the
firmware can find the functions by name. A label, a name followed by a colon,
stands for the address of what follows it. The only instruction of the
program, written three times, is \ins{ret}: return to the caller.

The cartridge is built and run with the commands met in
Chapter~\ref{ch:setup} (Listing~\ref{lst:fc:build}); the cartridge file is
byte for byte the same for QEMU, for the board and for \prgqt, where the
address is \code{http://localhost:8080}.

\shfile[caption={Building and running a cartridge.},label=lst:fc:build]{code/l02/c04.sh}

\begin{prgcheck}[Run it]
The expected result is a black screen. The achievement is real all the same:
the framework has found the three functions, has called them frame after frame,
and has survived.
\end{prgcheck}

\subsection{Step 2: say hello}

To put something on the screen the program calls two functions of the
framework: \code{prg32\_gfx\_clear}, which fills the screen with a colour, and
\code{prg32\_gfx\_text8}, which writes a string with a font of 8 by 8 pixels
(Listing~\ref{lst:fc:step2}, Figure~\ref{fig:fc:hello}). At the top of the
file two \code{.equ} directives give names to two constants:
\code{.equ BLACK, 0x0000} and \code{.equ WHITE, 0xffff}.

\asmfile[caption={Step 2: the new \code{first\_draw}.},label=lst:fc:step2]{code/l02/c05.S}

\begin{figure}[htb]
\centering
\includegraphics[width=0.55\textwidth]{lab1_step2_hello.png}
\caption{Step 2, built and run on \prg.}\label{fig:fc:hello}
\end{figure}

The rule for passing arguments is simple: the first argument goes in
\reg{a0}, the second in \reg{a1}, and so on. \ins{li} (\emph{load immediate})
puts a number in a register; \ins{la} (\emph{load address}) puts there the
address of a label. The string lives in the section \code{.rodata},
\emph{read-only data}, and the directive \code{.asciz} stores its characters
followed by a byte of value zero that marks its end.

\subsection{Why save \reg{ra}? The stack frame}

The first and the last three lines of \code{first\_draw} are new, and they are
the price of calling another function. The instruction \ins{call} stores in
\reg{ra} the address to come back to. But \code{first\_draw} was itself
called by the firmware, and \reg{ra} already held the address of
\emph{its} return: the first \ins{call} would overwrite it, and the final
\ins{ret} would jump to the wrong place.

A function that calls another one therefore saves \reg{ra} first. It lowers
\reg{sp} by 16 to make room on the \term{stack}, a region of memory that each
function may use for its own temporary needs, and stores \reg{ra} at
\code{12(sp)}, that is, 12 bytes above the address in \reg{sp}
(Figure~\ref{fig:fc:stack}). These lines are the \term{prologue} of the
function. After the last call it undoes everything in reverse order, in the
\term{epilogue}: it loads \reg{ra} back, raises \reg{sp} by 16, and returns.

\begin{figure}[htb]
\centering
\begin{tikzpicture}[x=1cm,y=0.72cm,font=\ttfamily\small]
  \foreach \i/\t/\o/\f in {0/caller's data/16(sp)/prgrule!60, 1/ra/12(sp)/prgamber!70, 2/free/{ 8(sp)}/prgpanel, 3/free/{ 4(sp)}/prgpanel, 4/free/{ 0(sp)}/prgpanel}{
    \draw[prgsteel,fill=\f] (0,-\i) rectangle (3.6,-\i-1);
    \node at (1.8,-\i-0.5) {\t};
    \node[anchor=west,text=prgcomment] at (3.8,-\i-0.5) {\o};
  }
  \draw[->,very thick,prgred] (-1.2,-4.5) -- (-0.05,-4.5) node[pos=0,left,font=\ttfamily\bfseries] {sp};
  \node[lbl,text=prgcomment] at (1.8,0.4) {higher addresses};
  \node[lbl,text width=7cm,align=center] at (1.8,-5.7) {The stack grows downward: subtract from \texttt{sp} to allocate, add to release.};
\end{tikzpicture}
\caption{Stack memory after the prologue of \code{first\_draw}.}
\label{fig:fc:stack}
\end{figure}

\begin{prgwarn}[The classic first bug]
Forgetting the pair \code{sw ra}/\code{lw ra} is the classic first bug: the
function returns to the wrong place, and the program crashes or freezes long
after the line that is at fault.
\end{prgwarn}

The convention requires \reg{sp} to stay aligned to a multiple of 16 bytes
across calls, which is why the frame is 16 bytes although only 4 are used.
The stack, with push and pop, is treated formally in
Chapter~\ref{ch:processor}.

\subsection{Step 3: a value that survives between frames}

Registers do not keep their values between one call of \code{update} and the
next: the firmware uses them in the meantime. Whatever must last has to live
in memory (Listing~\ref{lst:fc:step3}).

\asmfile[caption={Step 3: the new \code{first\_update} and its variable.},label=lst:fc:step3]{code/l02/c06.S}

\code{counter} is a label: it names a location in memory, in the section
\code{.data}, which holds one 32-bit word, initially zero. \ins{la} puts the
\emph{address} of that location in \reg{t0}. \ins{lw} reads the \emph{value}
at that address into \reg{t1}; \ins{sw} writes it back. Because the processor
computes only in registers, every change to a variable has the same three
movements: \emph{load, modify, store}. The function is a \term{leaf}: it
calls nothing, so it needs no stack frame.

\begin{prgconcept}[Address or value? The crux of memory]
A label names a \emph{location}. \ins{la} gives the location; \ins{lw} gives
the contents (Figure~\ref{fig:fc:addr}). This is how every variable of every
program ultimately works. Using an address where a value belongs, or the
reverse, is one of the most common beginner's errors: comment every register
as ADDRESS or VALUE.
\end{prgconcept}

\begin{figure}[htb]
\centering
\begin{tikzpicture}[font=\ttfamily\small]
  \node[lbl,anchor=west,font=\sffamily\small\bfseries,text=prgsteel] at (0,1.0) {Registers (inside the processor)};
  \node[anchor=east] at (0.6,0.2) {t0};
  \node[draw=prgaccent,fill=prgaccent!20,minimum width=42mm,minimum height=7mm,anchor=west] (t0) at (0.8,0.2) {address of counter};
  \node[anchor=east] at (0.6,-1.0) {t1};
  \node[draw=prgamber,fill=prgamber!25,minimum width=42mm,minimum height=7mm,anchor=west] (t1) at (0.8,-1.0) {41, then 42};
  \node[lbl,anchor=west,font=\sffamily\small\bfseries,text=prgsteel] at (8.2,1.0) {Memory (\texttt{.data} section)};
  \node[draw=prgsteel,fill=prgpanel,minimum width=38mm,minimum height=7mm,anchor=west] (m0) at (8.2,0.2) {...};
  \node[draw=prgsteel,fill=prgamber!15,minimum width=38mm,minimum height=7mm,anchor=west] (m1) at (8.2,-0.5) {counter: 41};
  \node[draw=prgsteel,fill=prgpanel,minimum width=38mm,minimum height=7mm,anchor=west] (m2) at (8.2,-1.2) {...};
  \draw[->,very thick,prgaccent] (t0.east) -- node[above,sloped,lbl,text=prgaccent]{\emph{t0 points here}} ([yshift=1.5mm]m1.west);
  \draw[->,very thick,prgamber] ([yshift=-1.5mm]m1.west) -- node[below,sloped,lbl,text=prgamber]{\emph{lw copies the contents}} (t1.east);
\end{tikzpicture}
\begin{asm}[numbers=none]
la   t0, counter   # t0 = WHERE the value lives
lw   t1, 0(t0)     # t1 = WHAT is stored there
addi t1, t1, 1     # change the copy in the register
sw   t1, 0(t0)     # memory still says 41 until this line
\end{asm}
\caption{The register \reg{t0} holds an address; \reg{t1} holds a value.}
\label{fig:fc:addr}
\end{figure}

Nothing visible changes after this step: the counter ticks, but nobody draws
it. How could one prove that it is ticking? With the debugger of \prgqt, or
with the next step.

\subsection{Step 4: make the counter visible}

Eight lines added to \code{first\_draw}, after the call that writes the text,
draw a green bar whose width is the counter (Listing~\ref{lst:fc:step4},
Figure~\ref{fig:fc:bar}). The constant \code{GREEN} is defined at the top of
the file by \code{.equ GREEN, 0x07e0}.

\asmfile[caption={Step 4: the lines added to \code{first\_draw}.},label=lst:fc:step4]{code/l02/c07.S}

\begin{figure}[htb]
\centering
\includegraphics[width=0.55\textwidth]{lab1_step4_bar.png}
\caption{Step 4: the bar grows from 0 to 255 pixels, then starts again.}
\label{fig:fc:bar}
\end{figure}

The function \code{prg32\_gfx\_rect} takes five arguments: $x$, $y$, width,
height and colour. The width is computed by \ins{andi} with 255, which keeps
only the low eight bits of the counter: the bar grows from 0 to 255 pixels and
then starts again. Three future chapters hide in these eight lines: the
equality $255 = \hex{FF}$ (positional notation), \ins{andi} as a bit mask
(Boolean algebra), and \hex{07E0} as a colour (the encoding of information).

The complete program, \fname{code/l02.S}, is Listing~\ref{lst:lab:l02} in
Section~\ref{sec:lab:l02}, with the commands that build it and run it.

% ----------------------------------------------------------------------------
\section{Reading what you wrote}

\subsection{The instructions just used}

Table~\ref{tab:fc:ins} orders by family the instructions met so far. In all
of them the result goes in the \emph{first} register written. The base
integer instruction set of \riscv, RV32I, has about forty instructions: a
quarter of them have already been used.

\begin{table}[htb]
\centering\small
\caption{The instructions used in this chapter.}\label{tab:fc:ins}
\begin{tabularx}{\textwidth}{@{}llLl@{}}
\toprule
\textbf{Family} & \textbf{Instruction} & \textbf{What it does} & \textbf{Example}\\
\midrule
Arithmetic & \ins{add}, \ins{sub} & combine two registers & \code{add t0, t1, t2}\\
Logic & \ins{and}, \ins{or}, \ins{xor} & bitwise operation on two registers & \code{and t0, t1, t2}\\
Immediate & \ins{addi}, \ins{andi}, \ins{ori}, \ins{xori} & one register and one small constant & \code{addi t1, t1, 1}\\
Memory & \ins{lw}, \ins{sw} & load or store a 32-bit word & \code{lw t1, 0(t0)}\\
Control & \ins{call}, \ins{ret} & go to a function and come back & \code{call prg32\_gfx\_rect}\\
\bottomrule
\end{tabularx}
\end{table}

\subsection{Pseudo-instructions}

Several of the mnemonics used are not instructions of the processor at all.
They are \term{pseudo-instructions}: shortcuts that the assembler expands into
one or two real instructions (Table~\ref{tab:fc:pseudo}). A
pseudo-instruction is not a separate hardware operation. The register
\reg{zero} is what makes many of them possible: a small instruction set,
reused with care.

\begin{table}[htb]
\centering\small
\caption{Pseudo-instructions and their expansions.}\label{tab:fc:pseudo}
\begin{tabular}{@{}lll@{}}
\toprule
\textbf{You write} & \textbf{The assembler emits} & \textbf{Meaning}\\
\midrule
\code{li t0, 42} & \code{addi t0, zero, 42} & load a constant\\
\code{mv t0, t1} & \code{addi t0, t1, 0} & copy a register\\
\code{not t0, t1} & \code{xori t0, t1, -1} & bitwise complement\\
\code{la t0, counter} & \code{auipc t0, \ldots} + \code{addi t0, t0, \ldots} & load the address of a label\\
\code{ret} & \code{jalr zero, 0(ra)} & jump to the address in \reg{ra}\\
\code{beqz t0, skip} & \code{beq t0, zero, skip} & branch if \reg{t0} is zero\\
\bottomrule
\end{tabular}
\end{table}

\subsection{One instruction is one 32-bit number}

The processor does not read text. The assembler translates each line into one
32-bit word, and it is that word that is stored in the cartridge and executed.
Figure~\ref{fig:fc:encoding} dissects \code{addi t1, t1, 1}: twelve bits hold
the constant, five bits name the source register (\reg{t1} is \reg{x6}), three
bits select the operation, five bits name the destination register and seven
bits, the \term{opcode}, say what kind of instruction this is. Written in
hexadecimal the thirty-two bits are \hex{00130313}.

\begin{figure}[htb]
\centering
\begin{bytefield}[bitwidth=0.029\textwidth,bitheight=8mm]{32}
\bitheader[endianness=big]{0,6,7,11,12,14,15,19,20,31}\\
\bitbox{12}[bgcolor=prgamber!30]{\texttt{000000000001}} & \bitbox{5}[bgcolor=prgaccent!20]{\texttt{00110}} & \bitbox{3}[bgcolor=prgviolet!20]{\texttt{000}} & \bitbox{5}[bgcolor=prggreen!20]{\texttt{00110}} & \bitbox{7}[bgcolor=prgsteel!20]{\texttt{0010011}}\\
\bitbox[]{12}{\footnotesize immediate\\the constant 1} & \bitbox[]{5}{\footnotesize rs1\\\texttt{t1}=\texttt{x6}} & \bitbox[]{3}{\footnotesize funct3\\add} & \bitbox[]{5}{\footnotesize rd\\\texttt{t1}=\texttt{x6}} & \bitbox[]{7}{\footnotesize opcode\\reg.-imm.}
\end{bytefield}
\caption{\code{addi t1, t1, 1} is the number \hex{00130313}.}
\label{fig:fc:encoding}
\end{figure}

\begin{prgconcept}
Assembly is a readable spelling of machine code. An instruction is data: the
processor reads its fields with logic circuits, which will be designed in the
chapters on logic and on the processor.
\end{prgconcept}

\subsection{PRG32-QT: a microscope for registers}

\prgqt runs the same cartridge on the desktop one instruction at a time, stops
between any two instructions and shows every register.

\begin{prgtry}[Step through \code{first\_update}]
Load step~3 in debug mode, place a breakpoint on \code{first\_update} and
press \key{Step} four times. Before each press, say what \reg{t0} and
\reg{t1} will contain; after it, check.
\end{prgtry}

\subsection{Trace it by hand}

The same exercise can, and should, be done on paper.
Table~\ref{tab:fc:trace} follows \code{first\_update} through a few frames.
At frame 256 the counter keeps growing, but its low eight bits return to zero,
because $256 = 1\,0000\,0000_2$: the bar starts again. Tracing by hand is the
most valuable debugging skill taught in this book.

\begin{table}[htb]
\centering\small
\caption{Three frames, and one more, of \code{first\_update}.}\label{tab:fc:trace}
\begin{tabular}{@{}ccccc@{}}
\toprule
\textbf{Frame} & \textbf{\ins{lw}: \reg{t1} becomes} & \textbf{\ins{addi}: \reg{t1} becomes} & \textbf{\ins{sw}: counter becomes} & \textbf{Bar width}\\
\midrule
1 & 0 & 1 & 1 & 1\\
2 & 1 & 2 & 2 & 2\\
256 & 255 & 256 & 256 & 0\\
257 & 256 & 257 & 257 & 1\\
\bottomrule
\end{tabular}
\end{table}

% ----------------------------------------------------------------------------
\section{What comes next}

The program of this chapter has raised more questions than it has answered.
Table~\ref{tab:fc:questions} is the forward map of the book: each later
chapter opens by pointing at one of its rows.

\begin{table}[htb]
\centering\small
\caption{Questions that the first program has already raised.}\label{tab:fc:questions}
\begin{tabularx}{\textwidth}{@{}P{0.29\textwidth}LP{0.2\textwidth}@{}}
\toprule
\textbf{You saw} & \textbf{The question} & \textbf{Answered in}\\
\midrule
\hex{07e0}, \hex{ffff}, 255 & Why do programmers write numbers in hexadecimal? & Chapters \ref{ch:information}, \ref{ch:integers}\\
a position below zero shown as \hex{FFFFFFFF} & How does a register hold a negative number? & Chapter \ref{ch:integers}\\
a ball that moves less than a pixel per frame & How are fractions represented without a floating-point unit? & Chapter \ref{ch:reals}\\
\code{andi a2, t1, 255} & What does AND do to thirty-two bits at once? & Chapters \ref{ch:boole1}, \ref{ch:boole2}\\
\code{addi t1, t1, 1} & Which circuit actually adds two numbers? & Chapter \ref{ch:boole2}\\
a pressed button is a 1 & What is a bit made of, physically? & Chapter \ref{ch:electronics}\\
\code{counter} survives between frames & How does a circuit remember a bit? & Chapters \ref{ch:sequential}, \ref{ch:memory}\\
\hex{00130313} & How does the processor decode and execute it? & Chapters \ref{ch:processor}, \ref{ch:instructions}\\
three functions and a rectangle & Can a whole game be made out of this? & \vermone: Chapters \ref{ch:verm1p1}, \ref{ch:verm1p2}, \ref{ch:verm1p3}, \ref{ch:verm1p4}\\
\bottomrule
\end{tabularx}
\end{table}

% ----------------------------------------------------------------------------
% Generated by tools/textbook/export_listings.js from decks/plan.js and code/l02.S: do not edit.
\section{The program of this chapter}\label{sec:lab:l02}

The listings of this chapter are parts of one program, \fname{code/l02.S} in the
repository of the course, printed in full in Listing~\ref{lst:lab:l02}. The
automatic tests of the course material build it and run it on the headless
runner of \prgqt{}, with scripted input, and check the frame;
Figure~\ref{fig:lab:l02} is the frame that the test wrote.

\begin{description}[style=nextline,leftmargin=1.5em]
\item[On the screen.] HELLO PRG32 and, under it, a green bar that grows by one pixel per frame and starts again at 256.
\item[What to press.] Nothing: the program runs by itself.
\end{description}

\begin{figure}[htb]
\centering
\includegraphics[width=0.5\textwidth]{lab_l02.png}
\caption{\fname{l02.S} on \prgqt{}.}\label{fig:lab:l02}
\end{figure}

\asmfile[caption={The complete program of the chapter, \fname{code/l02.S}.},label=lst:lab:l02]{code/l02/program.S}

To run it on \prgqt{} the commands are those of Listing~\ref{lst:lab:l02:run};
Section~\ref{sec:prg:testing} gives the same steps for QEMU and for the board.

\begin{shell}[caption={Building \fname{l02.S} and running it on \prgqt{}.},label=lst:lab:l02:run]
# 1. start PRG32-QT on this computer (it listens on port 8080)

# 2. in the PRG32 folder, with ESP-IDF active: build the cartridge
python3 -m prg32 cartridge build ~/arc12/code/l02.S --portable \
    --entry-prefix first --name l02 --out l02.prg32

# 3. send it to PRG32-QT and start it
python3 -m prg32 esp32c6 upload-and-run l02.prg32 \
    --url http://localhost:8080 --slot cart0

# 4. look, press, compare with the previous slide
\end{shell}


% ----------------------------------------------------------------------------
\section{Check your understanding}

\begin{quiz}{After \code{la t0, counter} and \code{lw t1, 0(t0)}, what does \reg{t0} hold?}
\item The value of the counter
\item The address of the counter
\item Zero, because \ins{la} clears the register
\item The return address
\end{quiz}

\begin{quiz}{A function that calls \code{prg32\_gfx\_clear} must save one register before the call. Which one?}
\item \reg{a0}, because it carries the argument
\item \reg{sp}, because the stack moves
\item \reg{ra}, because \ins{call} overwrites it
\item \reg{t0}, because it is temporary
\end{quiz}

\begin{quiz}{\reg{t1} holds 300. What does \code{andi a2, t1, 255} leave in \reg{a2}?}
\item 300
\item 255
\item 45
\item 44
\end{quiz}

\begin{quiz}{What number is \hex{1F}?}
\item the decimal number 1F
\item 31: one sixteen plus fifteen
\item the register t1F
\item a comment
\end{quiz}

\section{Exercises}

\begin{exercise}[on \prg]
\begin{enumerate}[nosep]
\item Change the title string and its colour to yellow (\hex{ffe0}); run it again.
\item Make the bar grow twice as fast, then four times as fast, changing one
  instruction only.
\item Make the bar shrink instead of grow. What happens when the counter
  passes zero?
\item In \prgqt, single-step \code{first\_update} and write down \reg{t0} and
  \reg{t1} after each instruction.
\item Draw a second bar, in red, whose width is the counter modulo 64.
\end{enumerate}
\expected{In the third task the counter goes from 0 to \hex{FFFFFFFF}: with
\code{andi 255} the bar jumps to 255 and then shrinks. Keep the observation:
Chapter~\ref{ch:integers} explains it.}
\end{exercise}

\section{Glossary}

\begin{glossarytable}
register & registro\\
load / store & caricare / memorizzare\\
address & indirizzo\\
immediate & immediato (costante)\\
branch & salto condizionato\\
label & etichetta\\
stack & pila (stack)\\
stack pointer & puntatore alla pila\\
return address & indirizzo di ritorno\\
calling convention & convenzione di chiamata\\
machine code & codice macchina\\
cartridge & cartuccia\\
\end{glossarytable}

\section{Summary and reading}

A \riscv processor is thirty-two registers and a small set of instructions
that compute, move and jump. A \prg game is three functions; arguments travel
in \reg{a0} to \reg{a7}; a function that calls another saves \reg{ra}.
Variables live in memory: \ins{la} for the address, \ins{lw} and \ins{sw} for
the value. An instruction is a 32-bit number; assembly is its readable
spelling.

\paragraph{Reading.} \emph{Playing with RISC-V}, chapter~5, ``First Contact''
\citep{prg32textbook}; \citet{patterson2020}, chapter~2; the tutorial
\fname{docs/learn/tutorial\_ascii\_game.md} of the \prg repository.
